% sup1_05_sbp_induction: Supplement part 1, steps S4 (the summation by parts (2.8)-(2.9) of Ba3) and S5 (the induction (2.16)).

\subsection{S4: the summation by parts \cite[(2.8)--(2.9)]{Ba3} at a vertex of \cite[(2.4)]{Ba3}}\label{sup1:ssec:S4}

The graph \cite[(2.4)]{Ba3} is two vertices \cite[(1.8)]{Ba3} with $n+n'=1$, joined by one scalar line carrying the derivative at both ends, with two external scalar and two external vector legs; its degree is $0$. \cite[p.~425]{Ba3} replaces such graphs by graphs of positive degree ``by integration by parts'' at one of the two vertices, \cite[(2.8)]{Ba3}, applied ``for all graphs $G_1,G_2^{(2)},G_3^{(3)},\dots$'' of a fixed ordering, so that $G'$ becomes a sum of graphs $\{G'^*\}$ whose components $G_1,\dots,G_m=G'^*$ along the ordering have positive degrees. Two things are to be proved: the identity, as an identity of multilinear forms in the legs, and the degree statement, which is the hypothesis of S5.

\begin{lemma}[the identity]\label{sup1:lem:S41}
Let $X\subset T_\eta$, $\phi_1,\phi_2\colon X\to\R^N$, and $h_\mu$ real functions with $\supp h_\mu\subset\{x:\dist_\infty(x,X^c)\ge2\eta\}$; $\tilde B$ arbitrary and real. Put $D^{\leftarrow}_\mu\phi(x)\defeq U(\tilde B_{x-\eta e_\mu,\mu})^{-1}(D_\mu\phi)(x-\eta e_\mu)=\eta^{-1}(\phi(x)-U^{-1}\phi(x-\eta e_\mu))$ and $d^*_\mu h(x)\defeq\eta^{-1}(h(x-\eta e_\mu)-h(x))$. Then
\[
(-e)\sum_\mu\sum_x\eta^d[(D_\mu\phi_1)(x)\cdot q\phi_2(x)]h_\mu(x)=(-e)\sum_x\eta^d[\phi_1(x)\cdot q\phi_2(x)]\Big(\sum_\mu d^*_\mu h_\mu\Big)(x)+e\sum_x\sum_\mu\eta^d[\phi_1(x)\cdot q(D^{\leftarrow}_\mu\phi_2)(x)]h_\mu(x-\eta e_\mu).
\]
With $h_\mu=gA'_\mu$, $d^*_\mu(gA'_\mu)(x)=g(x)(d^*_\mu A'_\mu)(x)+(d^*_\mu g)(x)A'_\mu(x-\eta e_\mu)$ gives the first two terms of \cite[(2.8)]{Ba3} exactly; the third printed term carries the derivative at the bond $(x-\eta e_\mu,x)$ transported to $x$, and its sign is as above.
\end{lemma}

\begin{proof}
For $F(x)\defeq q\phi_2(x)h_\mu(x)$, supported where both bonds at $x$ in direction $\mu$ lie in $X$, $\sum_x\eta^d(D_\mu\phi_1)(x)\cdot F(x)=\sum_x\eta^d\phi_1(x)\cdot(D^*_\mu F)(x)$ with $(D^*_\mu F)(x)=\eta^{-1}(U^{-1}_{x-\eta e_\mu,\mu}F(x-\eta e_\mu)-F(x))$, no boundary term arising because $\supp F$ is at depth $\ge2\eta$. $U$ commutes with $q$, so $D^*_\mu(q\phi_2h_\mu)=qD^*_\mu(\phi_2h_\mu)$, and $D^*_\mu(\phi_2h_\mu)(x)=-(D^{\leftarrow}_\mu\phi_2)(x)h_\mu(x-\eta e_\mu)+\phi_2(x)(d^*_\mu h_\mu)(x)$. Multiply by $-e$ and sum over $\mu$.
\end{proof}

\begin{remark}\label{sup1:rem:S4sign}
With $D^\eta_{\tilde B,\mu}$ the forward covariant difference \cite[(1.7)]{Ba1} evaluated at $x$, the printed third term of \cite[(2.8)]{Ba3} has the opposite sign; it is exactly right with $D^{\eta*}_{\tilde B,\mu}$, the symbol of the first line of \cite[(2.8)]{Ba3}, in place of $D^\eta_{\tilde B,\mu}$, so the print most likely dropped an asterisk. This is harmless for Proposition~1, which takes absolute values of every factor \cite[p.~426]{Ba3}, the graphical rule \cite[(2.9)]{Ba3} recording only where the derivative sits.
\end{remark}

\paragraph{Which vertex, and why the identity holds graph by graph.} The vertex of \cite[(2.4)]{Ba3} is \cite[(1.8)]{Ba3} with $n+n'=1$: $-e(L^k\eps)\sum_b\eta^d[(D^\eta_{\tilde B}\phi')(b)\cdot q\phi'(b_-)](g_kA'_b)^n(\tilde A_b)^{n'}$; after the localisation of \cite[p.~420]{Ba3}, $g$ carries the partition-of-unity factor of the vector leg, so $h_\mu=gA'_\mu$ ($n=1$) or $h_\mu=\tilde A_\mu$ ($n'=1$). The supports are at depth $\ge2\eta$ in $\Omega$: $\supp g\subset\Omega_1$ at distance $\ge r(t)-2$ from $\Omega^c$, and $\dist(\supp\tilde A,\partial\Omega)>2r(t)$. $E(G')$ is a sum over pairings; for a fixed pairing the value is obtained by substituting into each vertex form a function for each leg (a slice $G^\eta_{(j)}(\cdot,y)$ on $\Omega$ for a scalar leg, a vector slice for an $A'$-leg, or an external field). The vertex form is multilinear in its legs and \cref{sup1:lem:S41} is an identity of multilinear forms, so it holds after substitution, graph by graph, for every fixed ordering and index set $J(\tilde l)$, even though $\phi'(x)\cdot q\phi'(x)=0$ as a polynomial (the two legs are distinct in each graph).

\begin{lemma}[{degrees after \cite[(2.8)]{Ba3}}]\label{sup1:lem:S42}
Use the degree count \cite[(2.1)--(2.2)]{Ba3}: $D_G(v)=(\text{number of factors }\eta\text{ in }v)+\frac{2-d}2(\text{legs of }v\text{ on internal lines of }G)-(\text{differentiations in }v\text{ acting on internal lines of }G)+(\text{contour credit})$, $D(G)=\sum_vD_G(v)-d$.
\textup{(a)} $D(\text{\cite[(2.4)]{Ba3}})=0$.
\textup{(b)} Let $v$ be a vertex of a \cite[(2.4)]{Ba3}-subgraph $G_0=\{l_0\}$ and apply \cref{sup1:lem:S41} at $v$. In each of the three terms the graph has the same vertices and lines; only the place of the derivative at $v$ changes, from $l_0$ to the $A'$-leg (or $\tilde A$-leg), to $g$, or to the $\phi_2$-leg. For every connected $H\ni l_0$: $D_{\mathrm{new}}(H)=D_{\mathrm{old}}(H)+1$ if the new place is not an internal line of $H$, and $=D_{\mathrm{old}}(H)$ otherwise; in particular $D_{\mathrm{new}}(G_0)=1$.
\textup{(c)} For a connected $H$ containing $v$ but not $l_0$, and containing the line of the new place: $D_{\mathrm{new}}(H)=D_{\mathrm{old}}(H)-1$.
\textup{(d)} Fix an ordering $\tilde l$ and call the component $G^{(i)}_i$ created by $l(i)$ \emph{fresh} if $l(i)$ joins two vertices not met by $l(1),\dots,l(i-1)$ and these with $l(i)$ form a \cite[(2.4)]{Ba3} graph. Apply \cref{sup1:lem:S41} at one vertex of each fresh component. If every component $G^{(\alpha)}_i$ of $G'$ along $\tilde l$ has $D>0$ except the fresh ones, then in each resulting $G'^*$ every component $G^{(\alpha)}_i$ along $\tilde l$ has $D>0$.
\end{lemma}

\begin{proof}
(a) Each vertex of \cite[(2.4)]{Ba3} has $\eta$-count $d$, one internal scalar leg and one derivative on it, so $D(v)=d+\frac{2-d}2-1=\frac d2$ and $D=2\cdot\frac d2-d=0$. (b) The three terms keep the $\eta$-count and the legs; only ``differentiations acting on internal lines'' changes: for $H\ni l_0$ the old derivative was on an internal line, the new one is on an internal line exactly when $H$ contains the $A'$-line or the $\phi_2$-line at $v$ (the $g$-term never is). (c) For $H\not\ni l_0$ the old derivative was not on an internal line of $H$ and the new one may be. (d) A vertex \cite[(1.8)]{Ba3} has exactly one derivative leg; a fresh component uses the derivative legs of both its vertices, so each vertex lies in at most one fresh component, the fresh components are vertex-disjoint and the transformations commute. If $v$ lies in the fresh component created by $l_0=l(i_0)$, then $v$ is met by no earlier line, so every component $G^{(\alpha)}_i$ with $i\ge i_0$ containing $v$ contains $l_0$, and none with $i<i_0$ contains $v$. By (b) every component containing a transformed vertex has $D_{\mathrm{new}}\ge D_{\mathrm{old}}>0$ or is fresh with $D_{\mathrm{new}}=1$; the others are unchanged.
\end{proof}

\begin{remark}\label{sup1:rem:S4scope}
\cite[p.~425]{Ba3} states ``the degrees of these graphs are $\ge D(G_0)$'' for a graph $G_0$ containing the \cite[(2.4)]{Ba3} subgraph, where it holds by (b), and claims positivity only for the subgraphs $G_1,\dots,G_m=G'^*$ along the ordering, which is (d). By (c), a connected subgraph through the transformed vertex avoiding $l_0$ can lose one unit, so $G'^*$ need not have every connected subgraph of positive degree; the induction of S5 uses only (d). The transformation depends on the ordering, so the sum over orderings of \cite[(2.7)]{Ba3} is taken first, as printed. The new vertices are admissible for the bounds of S3 and S5: a derivative on $A'$ is a derivative at one end of a vector slice (\cref{sup1:lem:S36}(a), one extra $s^{-1}$); $|d^*g|\le c$ since $g_k$ is bounded with its derivatives up to second order \cite[p.~412]{Ba3}; a derivative on $\phi_2$ is a derivative at one end of a scalar slice (\cref{sup1:lem:S35}) or of the external field (bounded by $\norm{h\Phi}_{1,\alpha}$); the shifts by $\eta e_\mu$ keep every point within $\eta$ of the vertex localisation; the $\eta$-count is unchanged, and no line carries two derivatives at one end, the $\phi_2$- and $A'$-legs of \cite[(1.8)]{Ba3} carrying none.
\end{remark}

\subsection{S5: the induction \cite[(2.16)]{Ba3}}\label{sup1:ssec:S5}

After \cite[(2.13)--(2.14)]{Ba3} the quantity to bound is, for a graph $G$ with lines $l$, slice indices $j\in J(\tilde l)$ and cubes $\Delta(v)$,
\begin{equation}\label{sup1:eq:214}
\tilde E(G(j),\{\Delta(v)\})=\sum_{\Delta(v)}\prod_vs_{j(v)}^{d+e_v}\prod_ls_{j_l}^{d_l}\prod_l\exp\big[-\delta\,s_{j_l}^{-1}\dist_\infty(\Delta(v_l),\Delta(v'_l))\big],
\end{equation}
with $\delta\defeq\delta_1/2$, $s_j=L^j\eta$, $j(v)\defeq\min_{l\ni v}j_l$, $e_v\ge0$ the explicit $\eta$-power of a vertex \cite[(1.8)--(1.11)]{Ba3} ($n+n'-1$, $n+n'-2$, \dots\ from \cite[p.~413]{Ba3}; $0$ for the other vertices), the bound $\eta^{e_v}\le s_{j(v)}^{e_v}$ being the ``proper power'' of \cite[(2.14)]{Ba3}, and $d_l\defeq-(d-2)-(\text{derivatives on }l)+(\text{contour legs of }l)$, each line having two legs. Put $D(H)\defeq\sum_{l\in H}d_l+\sum_{v\in V(H)}e_v+d(|V(H)|-1)$ for a connected subgraph $H$.

\begin{lemma}[dictionary]\label{sup1:lem:S50}
$D(H)$ is Ba{\l}aban's degree \cite[(2.2)]{Ba3}, $\sum_{v\in H}D_H(v)-d$, with the contour credit going to legs of lines internal to $H$ at the line's own scale (\cref{sup1:lem:S36}(b)). $D(H)$ is an integer, so $D(H)>0$ implies $D(H)\ge1$.
\end{lemma}

\begin{proof}
Sum \cite[(2.1)]{Ba3} over $v\in H$: each internal line contributes two legs, $-(d-2)$ in total, its derivatives and its contour legs; each vertex contributes $d$ (the $\eta^d$ of its sum) and $e_v$.
\end{proof}

\begin{theorem}[{\cite[(2.16)]{Ba3} for all $i$, and \cite[(2.15)]{Ba3}}]\label{sup1:thm:S51}
Fix an ordering $\tilde l=(l(1),\dots,l(m))$ of the $m$ lines of a connected graph $G$; for $j\in J(\tilde l)$, $j_{l(1)}\le\dots\le j_{l(m)}\le k-1$. Let $G_i\defeq\{l(1),\dots,l(i)\}$ with endpoints, $G^{(\alpha)}_i$ its components, $W_i\defeq\sum_\alpha D(G^{(\alpha)}_i)$, and assume $D(G^{(\alpha)}_i)>0$ for all $i,\alpha$, so $W_i\ge1$. At stage $i$ (Ba{\l}aban's $G/G_i$) every component of $G_i$ is shrunk to one vertex $u$ carrying only its cube factor $s_{j(u)}^d$, every vertex not yet met keeps $s_{j(v)}^{d+e_v}$, $j(u)$ is the minimum index of the remaining lines at $u$, $\Delta(u)$ has side $s_{j(u)}$, and the remaining lines keep $s_{j_l}^{d_l}$ with rates $\delta_i>0$; call the resulting expression $\tilde E_i(G/G_i)$. Then for every fixed $j_{l(i+1)}\le\dots\le j_{l(m)}$,
\begin{equation}\label{sup1:eq:216}
\sum_{j_{l(1)}\le\dots\le j_{l(i)}\le j_{l(i+1)}}\ \sum_{\Delta\text{'s of }G}\tilde E\ \le\ C_i\,s_{j_{l(i+1)}}^{W_i}\sum_{\Delta\text{'s of }G/G_i}\tilde E_i(G/G_i),
\end{equation}
which is \cite[(2.16)]{Ba3}: the single factor $(L^{j_{l(i+1)}}\eta)^{\sum_\alpha D(G^{(\alpha)}_i)}$ and a constant depending on the number of lines (hence on $\bar n$), on $\delta_1$ and on $L$. The constants are $C_0=1$, $C_{i+1}\defeq C_ie^{\delta_i'n_i}S_d(\delta_i')(1-L^{-1})^{-1}$, $\delta_{i+1}\defeq\delta_i'\defeq\delta_i/(n_i+1)$, $n_i$ the number of other remaining lines at the far endpoint of $l(i+1)$, $S_d(\delta)\defeq\sum_{n\in\Z^d}e^{-\delta(|n|_\infty-1)_+}$; at $L=5$, $(1-L^{-1})^{-1}=\frac54$. At $i=m$: $\sum_{j\in J(\tilde l)}\sum_\Delta\tilde E\le C_m5^d$, which is \cite[(2.15)]{Ba3}.
\end{theorem}

\begin{proof}
Induction on $i$; at $i=0$ the statement is the definition. Assume \eqref{sup1:eq:216}$_i$, let $l\defeq l(i+1)$, $j'\defeq j_l$, and $u,u'$ its endpoints in $G/G_i$.
(1) $l$ is the lowest remaining line, so $j(u)=j(u')=j'$ and $\Delta(u)$, $\Delta(u')$ have side $s_{j'}$.
(2) \emph{Case $u\ne u'$.} Split $\exp[-\delta_is_{j'}^{-1}\dist(\Delta(u),\Delta(u'))]$ into $n_i+1$ equal factors of rate $\delta'\defeq\delta_i/(n_i+1)$. For another line $l'$ at $u'$ with far end $u''$ and $j_{l'}\ge j'$: $s_{j'}^{-1}\ge s_{j_{l'}}^{-1}$ and $\dist(\Delta(u),\Delta(u''))\le\dist(\Delta(u),\Delta(u'))+\operatorname{diam}_\infty\Delta(u')+\dist(\Delta(u'),\Delta(u''))$ with $\operatorname{diam}_\infty\Delta(u')\le s_{j'}\le s_{j_{l'}}$, so one share times $l'$'s factor is $\le e^{\delta'}\exp[-\delta's_{j_{l'}}^{-1}\dist(\Delta(u),\Delta(u''))]$: the line $l'$ now ends at $u$ with rate $\delta'$. The last share sums $\Delta(u')$ over the paving of side $s_{j'}$: $\sum_{\Delta(u')}\exp[-\delta's_{j'}^{-1}\dist(\Delta(u),\Delta(u'))]\le S_d(\delta')$, scale-free since the cubes are labelled by $s_{j'}\Z^d$ and $\dist\ge s_{j'}(|n|_\infty-1)_+$. The factor of $u'$ stays ($s_{j'}^d$ for a shrunk component, $s_{j'}^{d+e_{u'}}$ for a vertex not yet met).
(3) \emph{Shrinking.} Merge $u,u'$ into $\hat u$ with $j(\hat u)$ the minimum index of the remaining lines at $\hat u$, $\ge j_{l(i+2)}\ge j'$ (if no line remains, $\hat u$ is the last vertex). Fix $\Delta(\hat u)$ of side $s_{j(\hat u)}$ in the nested paving and sum $\Delta(u)\subset\Delta(\hat u)$: every remaining line at $\hat u$ has $\exp[-\cdots\dist(\Delta(u),\cdot)]\le\exp[-\cdots\dist(\Delta(\hat u),\cdot)]$, independent of $\Delta(u)$; there are $(s_{j(\hat u)}/s_{j'})^d$ cubes $\Delta(u)$ in $\Delta(\hat u)$, which with $u$'s factor $s_{j'}^d$ (times $s_{j'}^{e_u}$ if $u$ was not yet met) gives $\hat u$'s cube factor $s_{j(\hat u)}^d$. What is left at scale $j'$ is $s_{j'}^{d_l+d+E}$, $E$ the sum of $e_v$ over the endpoints of $l$ not met before.
(4) \emph{Case $u=u'$} ($l$ closes a loop or is a tadpole). There is no $\Delta(u')$ sum; what is left at scale $j'$ is $s_{j'}^{d_l}$ if $u$ was met before and $s_{j'}^{d_l+e_u}$ for a tadpole at a vertex not met before; shrinking is as in (3).
(5) \emph{Degrees.} With the stage-$i$ factor $s_{j'}^{W_i}$, the power of $s_{j'}$ is $W_{i+1}$: for $u\ne u'$, $D(\text{new})=D(\text{comp }u)+D(\text{comp }u')+d_l+d$ with $D(\{v\})\defeq e_v$ for a vertex not yet met (the vertex count adds $|V|+|V'|-1$); for $u=u'$, $D(\text{new})=D(\text{old})+d_l$; all other components are unchanged.
(6) \emph{The scale sum.} The only dependence on $j'$ is now $s_{j'}^{W_{i+1}}$ with $W_{i+1}\ge1$ (the $j(\cdot)$ of the vertices involve only $l(i+2),\dots$, independent variables $\ge j'$), so $\sum_{j'\le j_{l(i+2)}}s_{j'}^{W_{i+1}}\le(1-L^{-1})^{-1}s_{j_{l(i+2)}}^{W_{i+1}}$; if no line remains, the sum over $j'\le k-1$ is $\le(1-L^{-1})^{-1}$. This is \eqref{sup1:eq:216}$_{i+1}$.
(7) \emph{The end.} At $i=m$ one vertex remains with its cube factor; its cube sum, over cubes of side $s$ meeting a region of sup-diameter $\le4$ at weight $s^d$, is $\le5^d$. The rates obey $\delta_m\ge\delta/\prod_i(n_i+1)>0$, the $n_i$ being bounded by the number of legs, a function of $\bar n$.
\end{proof}

\begin{remark}\label{sup1:rem:S5}
(i) The induction uses only $W_{i+1}\ge1$ at each stage, which is weaker than the positivity of each component that \cref{sup1:lem:S42}(d) supplies. (ii) Booking each component's degree at its own next scale, as a vertex weight $D$ at $j(u)\ge j_{l(i+1)}$, gives a larger right side (since $s_{j_{l(i+1)}}\le s_{j(u)}$ and $D>0$); the printed common factor $s_{j_{l(i+1)}}^{W_i}$ is the sharper bookkeeping, and \cite[(2.16)]{Ba3} is correct as stated. (iii) The theorem is stated for general line dimensions $d_l$ and vertex powers $e_v$, which covers the lines of \cite[Prop.~2.2]{Ba3} as far as they are of the slice type; the generalised graphs themselves are the next step (part II of this supplement). (iv) The hypothesis ``every component along the ordering has positive degree'' is exactly what S4 gives for the graphs $G'^*$; it is not the hypothesis that every connected subgraph has nonnegative degree, which the $G'^*$ need not satisfy (\cref{sup1:rem:S4scope}): e.g.\ two \cite[(1.8)]{Ba3} vertices $v,v'$ with $n=1$ joined by a scalar line $l'$ carrying both derivative legs, with $v$'s $A'$-leg contracted by a vector line $l_A$ to a third \cite[(1.8)]{Ba3} vertex, all other legs external, have $D(\{l'\})=0$, $D(\{l_A\})=2$, $D(\{l',l_A\})=2$; in the ordering $(l_A,l')$ the line $l'$ is not fresh, is not transformed, and $G'^*=G$ keeps the zero-degree connected subgraph $\{l'\}$, while the components along this ordering have degrees $2$, $2$ and the theorem gives $O(1)$.
\end{remark}
